\section{总结与延伸}

\subsection{核心结论}

本节把全篇压缩成可复用判断。读者可以把这些结论当作检查表，用来判断一个 AI 应用到底有没有生成系统、产品内生数据和平台重构潜力。

最后把 EP123 压缩成十条结论。

\begin{enumerate}
\item AI 产品经理没有被模型替代，而是进入模型能力定义和 context 设计的前台。
\item 压缩让离散任务形成连续结构，语言是低成本表达世界的重要智能载体。
\item 数据决定智能边界，算力决定逼近速度，算法决定路径效率。
\item 公域数据利好基座模型，领域数据利好大厂，产品内生数据利好应用创业。
\item 视频生成机会不在单个 generate 按钮，而在生产方法、recipe 和商业闭环。
\item 生成系统由 DSL、context、agent/workflow、模型和评测共同构成。
\item 应用公司和模型公司边界会模糊，但数据飞轮不同。
\item 推荐系统是库存分配，生成系统是按需生产，旧平台分发权会被削弱。
\item 内容经济会从单点内容和注意力，转向 recipe、IP 和信任。
\item 创作者不会简单消失，而会从内容生产者变成生产方法、品味和信任的组织者。
\end{enumerate}

\subsection{开放问题}

最后保留开放问题，因为本期很多判断还处在推演阶段。生成系统、产销平台、recipe 资产化和信任货币，都需要未来几年用产品和市场来验证。

\begin{itemize}
\item 生成系统会先在哪类内容上替代推荐系统？
\item 平台会自己升级为产销平台，还是新平台先出现？
\item 产品内生数据如何被保护，避免被基座模型公司吸收？
\item Recipe、IP 和信任能否形成可计价资产？
\item AI 生成内容的版权归属，如何在“消费即创作”的场景中定义？
\item 小团队远程工作室是否能长期承载复杂生成系统？
\end{itemize}

\subsection{拓展阅读}

\begin{itemize}
\item EP130 张月光访谈：流程设计到上下文设计、AI Native 产品和 Agent。
\item EP128 Manus 访谈：Agent 产品化、复杂性风险和 AI 像制造业。
\item EP134 数据综述：数据金字塔、Data Recipe 和数据定价。
\item EP139 Agent 综述：Agent 技术史、OpenClaw Moment 和社会辐射。
\item Let's Verify Step by Step、Language Modeling Is Compression 等材料：理解过程监督、压缩和泛化。
\end{itemize}

\begin{importantbox}{最后的判断}
EP123 最值得留下的是一个新平台想象：当生产系统足够强，平台不再只是分配库存，而是按用户意图生产内容。旧互联网的注意力货币不会消失，但更高价值的资产会转向 recipe、IP、品味和信任。
\end{importantbox}

\end{document}
